\section{Sistemas de 1° Ordem}

\begin{frame}{Polos, Zeros e Estabilidade}
  \scriptsize
  \begin{columns}[t]
    \column{0.4\textwidth}
    \textbf{Forma geral}
    \begin{equation*}
      G(s)=K \frac{(s-z_1)(s-z_2)\cdots(s-z_m)}{(s-p_1)(s-p_2)\cdots(s-p_n)}
    \end{equation*}
    \textbf{Definições}
    \begin{itemize}
      \item Polo ($p_i$): raiz do denominador
      \item Zero ($z_i$): raiz do numerador
    \end{itemize}
    \textbf{Interpretação}
    \begin{itemize}
      \item Polos $\Rightarrow$ modos naturais ($e^{st}$), dinâmica e estabilidade
      \item Zeros $\Rightarrow$ moldam a resposta transitória
    \end{itemize}
    \column{0.6\textwidth}
    \textbf{Estabilidade}
    \begin{itemize}
      \item Estável se todos os polos têm parte real negativa
    \end{itemize}
    \textbf{Plano-$s$ (polos e zeros)}
    \begin{center}
      \begin{tikzpicture}[scale=0.8]
        % eixos
        \draw[->] (-3,0) -- (3,0) node[right] {Re};
        \draw[->] (0,-2) -- (0,2) node[above] {Im};
        % polos (X)
        \draw (-1.5,0.8) node {$\times$};
        \draw (-1.5,-0.8) node {$\times$};
        % zeros (O)
        \draw (-0.5,0) node {$\circ$};
        \draw (1,0.5) node {$\circ$};
      \end{tikzpicture}
      \\
      $\times$ polos \quad $\circ$ zeros
    \end{center}
  \end{columns}
\end{frame}

\begin{frame}{Sistema de 1ª Ordem}
  \scriptsize
  \begin{columns}[t]
    \column{0.4\textwidth}
    \textbf{Função de transferência}
    \begin{equation*}
      G(s)=\frac{K}{\tau s+1}
    \end{equation*}
    \textbf{Parâmetros}
    \begin{itemize}
      \item $K$: ganho ($\frac{\Delta y}{\Delta u}$)
      \item $\tau$: constante de tempo
    \end{itemize}
    \textbf{Constante de tempo}
    \begin{itemize}
      \item $\tau \Rightarrow 63\%$
      \item $2\tau \Rightarrow 86\%$
      \item $3\tau \Rightarrow 95\%$
      \item $5\tau \Rightarrow 99\%$
    \end{itemize}
    \column{0.6\textwidth}
    \textbf{Respostas (sistema linear e invariante no tempo)}
    \vspace{0.2cm}
    \begin{tabular}{c|c c c}
      & impulso & degrau & rampa \\
      \hline
      $x(t)$ & $\delta$ & $u(t)$ & $t$ \\
      $y(t)$ & $\frac{K}{\tau}e^{-t/\tau}$
      & $K(1-e^{-t/\tau})$
      & $K \cdot t-K \cdot \tau(1-e^{-t/\tau})$
    \end{tabular}
    \vspace{0.3cm} \\
    \textbf{Relação}
    \begin{itemize}
      \item impulso $\xrightarrow{\int}$ degrau $\xrightarrow{\int}$ rampa
      \item rampa $\xrightarrow{d/dt}$ degrau $\xrightarrow{d/dt}$ impulso
    \end{itemize}
    \vspace{0.2cm}
    \textbf{Polos e zeros}
    \begin{itemize}
      \item Polo ($p$): $\Rightarrow s=-\frac{1}{\tau}$
    \end{itemize}
  \end{columns}
\end{frame}

\begin{frame}{Efeito da posição do zero}
  \begin{equation*}
    G(s) = K \cdot \frac{\left(-\frac{s}{z} + 1\right)}{\tau s + 1}
  \end{equation*}
  \begin{columns}
    \column{0.5\textwidth}
    \begin{itemize}
      \item $z \ll -\frac{1}{\tau}$: salto inicial pequeno
        \begin{itemize}
          \item Quanto mais a esquerda menor o salto
          \item Sistema tende ao padrão de 1ª ordem
        \end{itemize}
      \item Aproximando do polo:
        \begin{itemize}
          \item Resposta mais rápida
          \item Salto inicial ficando maior
        \end{itemize}
      \item $z = -\frac{1}{\tau}$ (igual ao polo):
        \begin{itemize}
          \item Cancelamento polo-zero
          \item Resposta torna-se uma reta
        \end{itemize}
    \end{itemize}
    \column{0.5\textwidth}
    \begin{itemize}
      \item Entre polo e origem:
        \begin{itemize}
          \item Surge overshoot
          \item Aumenta ao se aproximar de $0$
        \end{itemize}
      \item Zero positivo ($z > 0$): Resposta inversa:
        \begin{itemize}
          \item Saída começa no sentido oposto
          \item Efeito muito intenso
        \end{itemize}
      \item Aumentando $z$:
        \begin{itemize}
          \item Resposta inversa diminui
          \item Sistema tende ao padrão de 1ª ordem
        \end{itemize}
    \end{itemize}
  \end{columns}
\end{frame}
